\subsection{\texorpdfstring{例：\(L^{p}\) 空间中的线性表示}{例：L\^{}\{p\} 空间中的线性表示}}

同样设 G 为在局部紧空间 X 上连续左作用的局部紧群。令 \(\beta\) 为 X 上支集为 X 的正测度。假设存在连续函数 \(\chi > 0\)，定义在 \(G \times X\) 上，使得对每个 \(s \in G\)，

\[
\boldsymbol {\gamma} (s) \beta = \chi (s ^ {- 1}, \cdot) \cdot \beta
\]

（这特别意味着 \(\beta\) 在 \(\mathbf{G}\) 下拟不变。）于是，\(\chi\) 是乘子。因为对给定的 \(s, t\)（属于 \(\mathbf{G}\)），有

\[
\begin{array}{r l} \boldsymbol {\gamma} (s) \boldsymbol {\gamma} (t) \beta = \boldsymbol {\gamma} (s) \big (\chi (t ^ {- 1}, \cdot) \cdot \beta \big) & = \big (\boldsymbol {\gamma} (s) \chi (t ^ {- 1}, \cdot) \big) \cdot \big (\boldsymbol {\gamma} (s) \beta \big) \\ & = \big (\boldsymbol {\gamma} (s) \chi (t ^ {- 1}, \cdot) \big) \cdot \chi (s ^ {- 1}, \cdot) \cdot \beta , \\ \boldsymbol {\gamma} (s t) \beta & = \chi (t ^ {- 1} s ^ {- 1}, \cdot) \cdot \beta , \end{array}
\]

因此

\[
\chi (t ^ {- 1}, s ^ {- 1} x) \chi (s ^ {- 1}, x) = \chi (t ^ {- 1} s ^ {- 1}, x)
\]

在局部 \(\beta\)-几乎处处成立，因而处处成立，因为 \(\chi\) 连续且 \(\beta\) 的支集为 \(X\)。

令 \(p \in [1, +\infty[\)。对每个 \(f \in \mathcal{L}_{\mathbf{C}}^{p}(\mathbf{X}, \beta)\) 及每个 \(s \in \mathbf{G}\)，令 \(\gamma_{\chi,p}(s)f\) 为 \(\mathbf{X}\) 上由下式定义的函数：

\[
\left(\gamma_ {\chi , p} (s) f\right) (x) = \chi (s ^ {- 1}, x) ^ {1 / p} f (s ^ {- 1} x).
\]

有

\[
\begin{array}{c} \int^ {*} | \chi (s ^ {- 1}, x) ^ {1 / p} f (s ^ {- 1} x) | ^ {p} d \beta (x) = \int^ {*} | f (s ^ {- 1} x) | ^ {p} \chi (s ^ {- 1}, x) d \beta (x) \\ = \int | f (x) | ^ {p} d \beta (x), \end{array}
\]

因此 \(\gamma_{\chi,p}(s)f\in\mathcal{L}_{\mathbf{C}}^{p}(\mathrm{X},\beta)\)。可见 \(\gamma_{\chi,p}(s)\) 是 \(\mathcal{L}_{\mathbf{C}}^{p}(\mathrm{X},\beta)\) 的等距自同态，并通过取商定义 \(\mathrm{L}_{\mathbf{C}}^{p}(\mathrm{X},\beta)\) 的等距自同态，仍记为 \(\gamma_{\chi,p}(s)\)。另一方面，\(\chi^{1/p}\) 显然是乘子，因此根据 No.~3 中得到的结果，\(\gamma_{\chi,p}\) 是 G 在 \(\mathrm{L}_{\mathbf{C}}^{p}(\mathrm{X},\beta)\) 中的线性表示。

命题 8. --- \(\gamma_{\chi,p}\) 作为 G 在 \(\mathrm{L}_{\mathbf{C}}^{p}(\mathbf{X},\beta)\) 中的线性表示，是连续且等距的。

令 \(f \in \mathcal{K}(\mathrm{X})\)。当 \(s\) 趋于 \(s_0\)（在 \(\mathrm{G}\) 中）时，\(\boldsymbol{\gamma}_{\chi,p}(s)f\) 趋于 \(\boldsymbol{\gamma}_{\chi,p}(s_0)f\)（在 \(\mathcal{K}(\mathrm{X})\) 中），因而也在 \(\mathrm{L}_{\mathbf{C}}^p (\mathrm{X},\beta)\) 中趋于它。由于 \(\boldsymbol{\gamma}_{\chi,p}(s)\) 是等距的，应用 No.~1 的注 2 即得命题 8。

对于不假设 \(\chi\) 连续的情形，参见 §4, Exer. 13。

命题 9. --- 假设每个函数 \(\chi(s,\cdot)\) 都有界。则 \(\gamma_{\chi}\) 保持 \(\mathrm{L}_{\mathbf{C}}^{p}(\mathrm{X},\beta)\) 不变，并且 \(\gamma_{\chi}\) 作为 \(\mathbf{G}\) 在 \(\mathrm{L}_{\mathbf{C}}^{p}(\mathrm{X},\beta)\) 中的线性表示是连续的。

令 \(f\in \mathcal{L}_{\mathbf{C}}^{p}(\mathrm{X},\beta)\)。则

\[
\begin{array}{r l} \int^ {*} | \chi (s ^ {- 1}, x) f (s ^ {- 1} x) | ^ {p} d \beta (x) \\ & \leqslant \sup _ {x \in X} \chi (s ^ {- 1}, x) ^ {p - 1} \int^ {*} | f (s ^ {- 1} x) | ^ {p} \chi (s ^ {- 1}, x) d \beta (x) \\ & = \sup _ {x \in X} \chi (s ^ {- 1}, x) ^ {p - 1} \int | f (x) | ^ {p} d \beta (x), \end{array}
\]

因此 \(\pmb{\gamma}_{\chi}(s)f\in \mathcal{L}_{\mathbf{C}}^{p}(\mathrm{X},\beta)\)，并且

\[
\left\| \boldsymbol {\gamma} _ {\chi} (s) \right\| \leqslant \sup _ {x \in X} \chi (s ^ {- 1}, x) ^ {1 / q},\tag{5}
\]

其中 \(q\) 表示与 \(p\) 共轭的指数。若 \(f \in \mathcal{K}(\mathrm{X})\)，则 \(\boldsymbol{\gamma}_{\chi}(s)f\) 趋于 \(\boldsymbol{\gamma}_{\chi}(s_0)f\)；这种趋近在 \(\mathcal{K}(\mathrm{X})\) 中进行，因而也在 \(\mathcal{L}_{\mathbf{C}}^{p}(\mathrm{X},\beta)\) 中进行，当 \(s\) 趋于 \(s_0\) 时。因此表示 \(\boldsymbol{\gamma}_{\chi}\) 将 \(\mathrm{G}\) 表示在 \(\mathrm{L}_{\mathbf{C}}^{p}(\mathrm{X},\beta)\) 中是连续的（No.~1，命题 2）。

若 G 在 X 上右作用，则有类似于 Nos. 3、4、5 中性质的性质。

特别地，若把 G 看作通过左平移或右平移作用于自身，并取 \(\chi = 1\)，就得到 G 在 \(\mathcal{C}(\mathrm{G})\)、\(\mathcal{K}(\mathrm{G})\)、\(\overline{\mathcal{K}(\mathrm{G})}\)、\(\mathcal{C}'(\mathrm{G})\)、\(\mathcal{M}(\mathrm{G})\)、\(\mathcal{M}^{1}(\mathrm{G})\) 中的左、右正则表示。若取 \(\beta\) 为 G 上的左（分别地，右）Haar 测度，并取 \(\chi = 1\)，就得到 G 在 \(\mathrm{L}_{\mathbf{C}}^{p}(\mathrm{G}, \beta)\) 中的左（分别地，右）正则表示。

